\subsection{5. Warum die Herkunftslücke nicht durch längeres Rechnen
verschwindet}\label{warum-die-herkunftsluxfccke-nicht-durch-luxe4ngeres-rechnen-verschwindet}

Sei \(\Pi_x=(-1)^{N_{f,x}}\) die lokale Fermionparität. Die bisher
angegebenen bankinternen Paarumwandlungen, Zahloperationen,
zahlenerhaltenden Symmetrielifts sowie reine Bosonverbindungen
kommutieren mit jeder \(\Pi_x\). Jede endliche Komposition, lineare
Kombination, Adjungierung und jeder durch solche Operatoren realisierte
einzelne Messzweig tut das ebenfalls. Starke beschränkte Grenzwerte
bleiben im Kommutanten der Parität.

Der benötigte Link erfüllt dagegen \(\Pi_xB=-B\Pi_x\) und
\(\Pi_yB=-B\Pi_y\), während er die Gesamtparität erhält. \textbf{Er
liegt nicht in dieser bekannten lokalen geraden Algebra.} Auch adaptives
Wiederholen und Nachselektieren gerader Krauszweige kann die
Sektorgrenze nicht überwinden.

Die Voraussetzung „jeder realisierte Zweig ist gerade`` ist wesentlich.
Eine bloß paritätskovariante CP-Abbildung kann ungerade Krausoperatoren
besitzen. Der im Prüfer enthaltene Reset-Kanal auf einer Fermionmode ist
ein exaktes Gegenbeispiel zur unzulässigen stärkeren Behauptung.

Das beweist keinen universellen Unmöglichkeitssatz für den Compiler. Es
beweist, \textbf{welche neue Quelleneigenschaft gesucht werden muss}:
eine gerade Gesamtoperation, die bezüglich der zwei operational
definierten Teilbereiche ungerade ist, oder eine ursprüngliche
Überlappung, durch die die angenommene Zerlegung in unabhängig gerade
Banken falsch war. Eine andere bloße Clockphase oder weitere
Boson-Paarumwandlung reicht nicht.

\subsection{6. Was von der arithmetischen Untersuchung
trägt}\label{was-von-der-arithmetischen-untersuchung-truxe4gt}

\subsubsection{6.1 Eine einzelne endliche Bank ist nicht der volle
Logarithmusgenerator}\label{eine-einzelne-endliche-bank-ist-nicht-der-volle-logarithmusgenerator}

Quadratisches Ergänzen liefert mit der nativen Casimirnorm
\(\|\sum P_A^\dagger P_A\|\le480\)

\[\adjustbox{max width=\linewidth}{$\displaystyle 
 H_{\rm nat}\ge\tfrac\Delta2N_b-C,
 \qquad C=960g^2/\Delta.
$}\]

Die externe Rechnung verwendet teilweise die schwächere Dreiecksgrenze
\(C=7680g^2/\Delta\), die ebenfalls gültig ist. Über Minmax,
\textbf{nicht} eine allgemeine Operator-Monotonie der
Exponentialfunktion, folgen

\[\adjustbox{max width=\linewidth}{$\displaystyle 
 N_H(E)\le2^{64}\binom{\lfloor2(E+C)/\Delta\rfloor+60}{60},
 \qquad \operatorname{Tr}e^{-\beta H}
 \le\frac{2^{64}e^{\beta C}}{(1-e^{-\beta\Delta/2})^{60}}.
$}\]

Das Zustandswachstum ist polynomial, während ein vollständiges Spektrum
\(E_*\log n+E_{\rm off}\), \(E_*>0\), exponentiell viele Zustände unter
wachsender Energie verlangt. Die ursprüngliche Bank kann dieses
komplette Spektrum daher nicht energieerhaltend mit nur affiner
Skalenänderung tragen. Diese Schranke widerlegt nicht jeden
arithmetischen Untersektor, jede nichtlineare Umparametrisierung oder
jeden Nullstellenoperator.

Endlich viele solche Banken und endlich viele energetisch kontrollierte
Rotoren ändern den qualitativen Gegensatz nicht. Unendlich viele
identische Banken lösen ihn nicht automatisch: Bei gleichmäßig
beschränkten lokalen Anregungskosten divergiert bereits die globale
Wärmespur. Ein geeigneter relativer oder lokaler Spurbegriff wäre ein
zusätzlicher Gegenstand.

\subsubsection{6.2 Primzahlen können eine gewählte Dynamik organisieren,
erzwingen sie aber
nicht}\label{primzahlen-kuxf6nnen-eine-gewuxe4hlte-dynamik-organisieren-erzwingen-sie-aber-nicht}

Auf \(\ell^2(\mathbb N)\) kann man
\(H_{\rm ar}|n\rangle=\log n|n\rangle\) definieren. Die additive
Primfaktorbesetzung erklärt dann die Eulerstruktur. Der
Definitionsbereich lautet \(\sum_n(\log n)^2|\psi_n|^2<\infty\). Dass
quantenstatistische Systeme mit Zeta-Zustandssumme existieren, ist
bereits ein Ergebnis von
\href{https://alainconnes.org/wp-content/uploads/bostconnesscan.pdf}{Bost
und Connes}. Das ist ein relevanter Anschluss, nicht die fehlende native
Herleitung.

Der gelieferte Cocycle \(c(m,n)=(-1)^{v_2(m)v_3(n)}\) ist assoziativ,
weil Bewertungen unter Multiplikation additiv sind. Für
\(T_m|n\rangle=c(m,n)|mn\rangle\) gilt \(T_lT_m=c(l,m)T_{lm}\),
insbesondere \(T_2T_3=-T_3T_2\). Die unkontrollierten CP-Kanäle
verlieren diese globale Phase und erfüllen
\(\mathcal E_l\mathcal E_m=\mathcal E_{lm}\). In kohärent kontrollierten
Wegen bleibt die relative Phase messbar. Das ist eine exakte
Phasenstruktur, aber kein Beweis einer Quellenauswahl dieses Cocycle
oder von RH.

Strikt multiplikative Schritte \(m>1\) ergeben auch keine geschlossenen
Bahnen: \(n\mapsto mn>n\). Die Frequenz \(\log p\) eines Operators, eine
Bahnlänge \(r\log p\) und der Schwingungszeitraum \(2\pi/\log p\) müssen
auseinandergehalten werden. Hier hat der RH-Graph-Skill die Prüfung auf
Generator, Spur und Phase getrennt und eine bloße Analogieschließung
verhindert.

\subsubsection{6.3 Der RH-Zielsatz bleibt
unbewiesen}\label{der-rh-zielsatz-bleibt-unbewiesen}

Die vorgeschlagene Identität

\[\adjustbox{max width=\linewidth}{$\displaystyle 
 \frac{\Xi(z)}{\Xi(0)}=\det(I-z^2A^{-2}),\quad
 \Xi(z)=\xi(\tfrac12+iz),
$}\]

für einen strikt positiven selbstadjungierten Operator \(A\) mit
kompakter Inverser und \(A^{-2}\) von Spurklasse wäre tatsächlich
hinreichend: Die gesamte rechte Funktion hat ihre Nullstellen nur bei
reellen \(\pm\lambda_j(A)\), einschließlich Multiplizitäten. Aber dieser
Operator und die gesamte Funktionsidentität sind \textbf{nicht
konstruiert}. Ein Operator mit lediglich den Imaginärteilen der
Nullstellen wäre unzureichend.

Weder diese Prüfung noch die im externen Code bereits verwendete
Faktorisierung kleiner Testzahlen liefert einen neuen effizienten
Faktorisierungsalgorithmus, eine native Quantenoperation dafür oder eine
Lösung von P versus NP. Hylæan wurde in dieser Revision nicht neu
geprüft.

\subsection{7. Die lokale Polnäherung wird präziser, das relativistische
Feld bleibt
offen}\label{die-lokale-polnuxe4herung-wird-pruxe4ziser-das-relativistische-feld-bleibt-offen}

Die externe Resolventenabschätzung ist korrekt: Im Kreis mit Radius
\(.1\Delta\) um den \textbf{tatsächlichen}, unbekannten negativen Pol
\(-\epsilon\) gilt

\[\adjustbox{max width=\linewidth}{$\displaystyle 
 G(z)=\frac{Z}{z+\epsilon}+R(z),\qquad
 |R(z)|<.505213/\Delta.
$}\]

Der übrige Spektralträger ist mindestens \(.337373\Delta\) vom Pol
entfernt. Die relative Abweichung vom echten Polterm beträgt im Kreis
unter 5,75 \%. Die Abschätzung setzt keine Intervallmitte anstelle von
\(Z,\epsilon\) ein und ist keine globale Vernachlässigbarkeit der
Selbstenergie.

Aus v1.6.4 bleibt bestehen: Die exakte Zwei-Feld-Resolvente besitzt
einen nichtverschwindenden Rest \(\Sigma_2(z)\). Seine komplette
Funktion ist noch nicht bestimmt. Das neue kontrollierte Transportlemma
benötigt sie nicht vollständig, weil es den Rest durch eine
Spektrallücke kontrolliert.

Ebenso bleibt der Feldtypentest bestehen: Der native antisymmetrische
Paartensor zusammen mit zwei gleichhändigen Weylfeldern und einem
skalaren Vermittler ergibt den verschwindenden Kanal. Ein symmetrischer
Lorentz- Spinortensor oder eine zusätzliche gleichgeladene
zweidimensionale Antisymmetriemarke eröffnet einen nichtverschwindenden
Kanal, führt aber zusätzliche Struktur ein. Ein internes
Spin(10)-Spinorlabel ist kein Lorentz-Weylindex. Der neue Link behebt
diese Typfrage nicht. Die Konventionen des bisherigen Tests sind mit der
Primärübersicht von \href{https://arxiv.org/abs/0812.1594}{Dreiner,
Haber und Martin} verbunden; ein neuer vollständiger relativistischer
Adapter wurde hier nicht konstruiert.

Die spätere Feldtypenprüfung in \codeword{LATE\_AUDIT.md} zeigt außerdem
einen Lorentz-erlaubten Zwei-Ableitungs-Kandidaten für den symmetrischen
Spinortensor und widerlegt einen pauschalen (1,1)-Ausschluss: Der
Energie-Impuls-Tensor ist schon bei Dimension vier zulässig. Weder ein
gesunder nativer Feldadapter noch ein dynamischer Gravitonpol folgt
daraus.

\subsection{8. Während der Arbeit entdeckte Änderung: „Kommutant 7``
richtig
einordnen}\label{wuxe4hrend-der-arbeit-entdeckte-uxe4nderung-kommutant-7-richtig-einordnen}

Die strikte Originalquellenprüfung stoppte korrekt, als eine andere
Arbeit den früheren nativen Bericht ergänzte. Der alte Stand mit Hash
\codeword{1062e13fe0fb79cce015657a7b4d40a61e9639922ff95fa60c689c0068c38ffc}
blieb eingefroren. Der gesonderte geänderte Entwurf hat Hash
\codeword{ab745e74b1232b1ddf35dd0e0d0da4ace0176c797d769e5db8bedd0010bcba6a}.
Die in dieser Revision verwendeten Energie- und Polzertifikate änderten
sich bei dieser Prüfung nicht. Der Fehlerbeleg wurde aufbewahrt.

Der neue Entwurf listet drei dunkle Darstellungstypen mit Dimensionen
24.000, 11.200, 2.688; drei helle mit 2.880, 576, 320; und einen χ-Typ
mit 64. \textbf{Unter Voraussetzung dieser Zerlegung} lässt sich seine
Folgerung exakt nachprüfen:

\[\adjustbox{max width=\linewidth}{$\displaystyle 
 \mathcal H_{N=3}\simeq
 \bigoplus_{d\in D}V_d\ \oplus\
 \bigoplus_{b\in B}(\mathbb C^2\otimes V_b)\ \oplus V_\chi.
$}\]

Die Summe ist \(37888+2\cdot3776+64=45504\). Die drei hellen Typen
kommen im \textbf{vollen} Raum jeweils zweimal vor. Daraus folgen:

\begin{itemize}
\tightlist
\item
  Nur die Gruppenwirkung hat einen Kommutanten der Dimension
  \(3+3\cdot4+1=16\), nicht sieben.
\item
  Werden zusätzlich die nichtverschwindenden hellen Paarmischungen und
  \(N_b\) als Operationen zugelassen, erzeugen sie die vollständigen
  \(M_2\)-Multiplizitätsalgebren. Dann bleiben sieben zentrale Skalare.
\item
  Sieben ist \textbf{keine absolute Untergrenze für beliebige weitere
  Operationen}. Ein erlaubter Operator, der diese sieben Blöcke
  verbunden mischt, zwingt die sieben Skalarwerte zur Gleichheit und
  lässt nur einen übrig. Der endliche Inzidenzmatrix-Test hat Rang sechs
  und beweist dieses Gegenargument.
\item
  Die Zerlegung allein beweist weder operative Verfügbarkeit der
  kontinuierlichen Symmetrie noch einen Einzelfermion-Link. Selbst alle
  bankinternen Symmetrieoperationen erhalten die lokale Fermionparität.
\end{itemize}

Der interessante Wert sieben wird dadurch nicht verworfen. Korrigiert
werden die Aussagen „voller N=3-Raum multiplizitätsfrei``, „absolute
Grenze für jede Operation`` und eine unbewiesene Gleichsetzung von
Symmetrie und verfügbarem Instrument. Die genaue Darstellungszerlegung
selbst bleibt in diesem Zusatz eine externe Voraussetzung, kein frisch
reproduzierter Satz.

Der Debugging-Skill führte hier zur Versionsprüfung statt zum Ersetzen
eines Pins. Das neue Manifest trennt deshalb einen erfolgreichen
\textbf{eingefrorenen Quellenreplay} von
\codeword{originals\_unchanged=false}.

\subsection{9. T1-T8: was dieses Ergebnis beiträgt und nicht
beiträgt}\label{t1-t8-was-dieses-ergebnis-beitruxe4gt-und-nicht-beitruxe4gt}

Die Begriffe folgen der aktuellen offenen Problemliste und ihrer
Forschungszuordnung. Keine Akzeptanzmarkierung wurde heraufgesetzt.

{\def\LTcaptype{none} % do not increment counter
\begin{longtable}[]{@{}
  >{\raggedright\arraybackslash}p{(\linewidth - 4\tabcolsep) * \real{0.3333}}
  >{\raggedright\arraybackslash}p{(\linewidth - 4\tabcolsep) * \real{0.3333}}
  >{\raggedright\arraybackslash}p{(\linewidth - 4\tabcolsep) * \real{0.3333}}@{}}
\toprule\noalign{}
\begin{minipage}[b]{\linewidth}\raggedright
Tor
\end{minipage} & \begin{minipage}[b]{\linewidth}\raggedright
Fehlender Nachweis
\end{minipage} & \begin{minipage}[b]{\linewidth}\raggedright
Beitrag dieser Revision
\end{minipage} \\
\midrule\noalign{}
\endhead
\bottomrule\noalign{}
\endlastfoot
T1 & Ursprung von P1/P2, Dimension, Compiler- und Zustandswahl &
Bekannte gerade Operationsalgebra ist für Einzeltransport unzureichend;
ursprüngliche Auswahl weiter offen \\
T2 & Tatsächlich halbgeladenes, markiertes E8-Feld samt Renormierung,
Energie- und Adjungiertenkontrolle & Der Ladung-eins-Lochzustand ist
kontrolliert, aber nicht mit dem fehlenden Half-Charge-Feld
identifiziert \\
T3 & Ein gemeinsamer, aus TFPT ausgewählter lokaler unitärer 3+1D-Parent
& Vollständiger bedingter Zwei-Banken-Transfertest; kein ausgewählter
räumlicher Parent \\
T4 & Chirales SM-Maß, Anomalien/Index, gleichmäßige Spiegelentkopplung &
Feldtyp-Widerspruch weiter sichtbar; kein chirales Maß \\
T5 & Wechselwirkender Kontinuumsübergang, Lorentzverhalten, Clustering,
Confinement und Streuung & Endliche-Sektor-Fehlerkontrolle ist ein
Baustein, aber kein räumlicher Grenzübergang \\
T6 & Interne Herleitung aller Eichkopplungen und vollständiger
Neutrinotextur/-skala & Keine neue Herleitung; \(\tau,\kappa\) sind
deklarierte Testparameter \\
T7 & Masseloser dynamischer Spin zwei, beide Helizitäten, universelle
Kopplung im selben Parent & Offen; der neue pauschale Ausschluss eines
Dimension-vier-(1,1)-Tensors wurde korrigiert \\
T8 & Physische Präparation/Anfangszustand und ein gemeinsames
Quellfunktional aller Auslesungen & Anfänglicher Polfilter wird für
\textgreater89,7 \% unnötig; geladene Präparation und Messinstrument
weiter offen \\
\end{longtable}
}

\subsection{10. Nächste entscheidende Arbeiten mit klaren
Abbruchkriterien}\label{nuxe4chste-entscheidende-arbeiten-mit-klaren-abbruchkriterien}

\textbf{A. Die Quelle des geraden Gesamtlinks finden.} Nicht noch eine
riesige Kontrollmatrix bauen, sondern für jeden tatsächlich
ursprünglichen Generator \(O\) prüfen, ob \([O,\Pi_x]\ne0\) bei
erhaltener Gesamtparität möglich ist. Eine positive Antwort muss
Generator, Teilbereichsdefinition, Hilbertraum, Gaussreferenz und
Matrixelement liefern. Bleiben alle Generatoren lokal gerade, ist die
Suche nach dem Link durch deren weitere Komposition beendet. Dann muss
die ursprüngliche Teilraumdefinition oder der Quellenvertrag explizit
geändert werden; man darf keinen Hopterm als Ergebnis ausgeben.

\textbf{B. Die ursprüngliche Entnahme als gemeinsames Instrument
realisieren.} Eine Referenzmode könnte Ladung erhalten, während
\(f_r\Omega\) präpariert und ausgelesen wird. Zu prüfen sind ihre
Herkunft, Energiekosten, Erfolgshäufigkeit und Erhaltung der internen
Kohärenz. Der neue Satz entfernt bereits die zusätzliche Pflicht eines
perfekten anfänglichen Energiefilters. Er entfernt nicht die Pflicht,
ein geladenes Instrument auszuweisen.

\textbf{C. Den Feldtyp vor räumlichen Rechnungen entscheiden.} Entweder
folgt ein passender Lorentzträger wirklich aus derselben Quelle, oder
der skalare Weyl-Ansatz wird in seiner jetzigen Form verworfen. Eine
bloße zusätzliche Verdopplung ohne Herkunft ist ein Modellvorschlag,
keine Ableitung.

\textbf{D. Erst dann viele Banken.} Der nächste räumliche Test benötigt
dieselbe Quelle, einen bestimmten Graphen, Lokalisierung der Operationen
und mit dem Volumen verträgliche Fehler- und Energiebounds. Ein einziges
neues Linklemma liefert weder Raumdimension noch Lichtkegel, chirale
Materie oder Gravitation.

\textbf{E. Arithmetik getrennt scharf halten.} Eine echte Verbindung
müsste eine gemeinsame Quellabbildung mit Generator,
Zustands-/Spurstruktur, Phasen und Domänen liefern. Die volle signierte
Weilform oder die ganze Determinanten- identität bleibt das Beweisziel.
Der bloße Auftritt von Primzahlen erfüllt keines davon. Der globale
RH-Index ließ sich wegen einer fehlenden historischen Quelle und
veränderter Review-Quellen nicht vollständig aktualisieren; die gezielte
Originalquellenprüfung ist enger als ein frischer Gesamtindex. Es wurde
kein RH-Abschlusskandidat registriert.

\subsection{11. Reproduktion und ehrliche
Liefergrenze}\label{reproduktion-und-ehrliche-liefergrenze}

\codeword{replay.py --frozen-only} führt die fünf Prüfer normal und mit
\codeword{-OO} aus, vergleicht beide JSON-Ergebnisse und die zwei
externen Originalberichte, prüft alle eingefrorenen Quellen und
protokolliert Änderungen der Originalorte gesondert. Ein Fehler führt zu
FAIL. Der zuvor aufgetretene Live-Quellenfehler bleibt in
\codeword{source\_drift\_failure\_receipt.json} sichtbar.

Die abgesicherten Zahlen stehen vollständig rational in
\codeword{two\_bank\_transfer\_normal.json}; Dezimalwerte dienen nur der
Lesbarkeit. Die analytische Beweiskette steht in Abschnitt 4. Das Paket
umfasst die unveränderten externen Arbeiten, deren Programme, die native
v1.6.4-Abhängigkeit und alle neuen Berichte. NumPy, SymPy und für das
ältere native Paket dessen separat dokumentierte Voraussetzungen sind zu
unterscheiden.

Die versionierte ausführliche Markdown-Lieferung enthält zusätzlich den
vollständigen eingefrorenen nativen Herleitungsbericht v1.6.4 als
historischen Anhang. Daneben gibt es ein kurzes Update und eine
bildliche Erklärung. \textbf{Die großen TFPT-Haupt-PDFs und die Webseite
werden in dieser Revision nicht als aktualisiert ausgegeben. Es gab
keinen Commit oder Push.}

\begin{center}\rule{0.5\linewidth}{0.5pt}\end{center}

\section{Nachprüfung der beiden zuletzt eingegangenen
Untersuchungen}\label{nachpruxfcfung-der-beiden-zuletzt-eingegangenen-untersuchungen}

Integrierter Nachtrag zu v1.6.5 · 15. September 2026

Quellen: die Anlagen \codeword{8e110ae1-19e2-4d47-bfc8-ebe4fe9ab81b} und
\codeword{b928b72a-6669-4c67-a75d-124095791c6f}, unverändert
eingefroren. Zusätzlich wurden die vier genannten Programme
\codeword{t1\_fixed.py}, \codeword{stabilizer.py},
\codeword{t5\_twobank.py}, \codeword{t5\_varresp.py} vollständig gelesen
und archiviert. Eine eigene Nachrechnung prüft die kritischen
Folgerungen; sie ist \textbf{kein vollständiger Replay aller sechs
Sonden} des fremden Arbeitsordners.

